\section{Supplementary local execution audit}
\label{sec:execution}
The preceding travel-harness results observe plans. We now add a separate local
execution study with declared permission truth. It does not retroactively establish
the correctness of travel bookings or make the original policies comparable.
The purpose is to test whether log consistency and absence of effects can hide
different operational outcomes; action binding itself is established prior work.

\subsection{Adapters, permissions, and component baseline}
Three adapters perform distinct local operations: append one text line to a target
file; insert a transaction and update an account balance in SQLite; and enqueue a
structured message in a recipient-specific local outbox. No queued message is sent
and no real account, booking, or payment is changed. These are three adapter types,
not three independently sampled application domains.

Each adapter has four authored goals with different object and payload identifiers.
Index zero is development, while indices one through three form nine evaluation
base goals. Identifier separation does not establish a new external distribution.
Every episode resets its state. Permission records grant exact tool, target and
value tuples under a policy version and permit one use per action. Pending records
grant nothing. The same full goal/permission is supplied across guards; record IDs
remain outside model input.

We run the unmodified Microsoft AGT approval\_protocol component at revision
\texttt{47abed7446d4}, including its actual action binding, coordinator and in-memory
store~\citep{microsoft2026actionbound}. Its full revision, source hashes and MIT
licence are retained. The official component test module passes 31 tests under
import-only package scaffolding. This evaluates a component integration, not the
complete AGT framework. A trusted synthetic service approver replaces a real human
authentication flow. A declared fixed component clock controls expiry; wall-time
metadata from dataclass defaults is not a real grant chronology. Expiry behavior,
trusted approval presentation, total mediation, cross-process persistence and
distributed crash consistency are outside this experiment.

\subsection{Constructed faults and separate persisted-state measurement}
Six guards are allow all, deny all, name-only checking, a one-time exact precheck,
a local execution-bound implementation control, and native AGT. The local bound
control is not a proposed new defense. Ten states cover a clean grant, pending
approval, target/value changes before review, target/value changes after review,
sequential and two-thread duplicate calls, a policy-version change, and operation
failure. The evaluation contains 540 cells. These interventions are authored fault
fixtures; their frequency does not estimate natural failure prevalence.

The exact-precheck arm reuses its review Boolean after a change, while the bound
arms revalidate the actual action and current policy. A consumed approval is not
silently restored after operation failure. Threaded native calls rely on the
upstream store's consumption primitive; only adapter writes are serialized to avoid
unrelated storage corruption. No barrier or admission-interval trace establishes
overlap, so this fixture does not certify coverage of race schedules. The actual operation errors are filesystem directory
append errors or SQLite CHECK failures, rather than injected success labels.

A separately implemented verifier reloads file bytes, transaction rows, balances
and outbox contents. It compares fields directly with the original permission and
counts effects beyond its allowed use; it never uses a gate digest or claimed
success as authorization truth. It checks action-level permissions, not receipt
authenticity inferred from final effects. A second standalone checker additionally
derives permission truth from the frozen base-goal manifest and independently
recomputes native binding hashes. These programs have the same parent author;
implementation separation is not independent human review. Twenty-three finite
positive/negative controls include false success claims and corrupted persisted
state. Neither these controls nor the upstream tests are scientific efficacy trials.
\begin{table}[t]
\centering\small
\caption{Constructed local-effect audit on nine evaluation base goals. Each guard
has 90 cells. U counts unauthorized persisted effects; M counts cells with U and
positive action-name/log consistency. Clean utility uses nine granted goals.
These are authored contract cases, with duplicate effects counted separately.}
\label{tab:execution-contract}
\begin{tabular}{lrrrr}\toprule
Guard & Effects & U & M & Clean utility\\\midrule
Allow all & 99 & 72 & 72 & 9/9 \\
Deny all & 0 & 0 & 0 & 0/9 \\
Name only & 90 & 63 & 63 & 9/9 \\
Exact precheck & 72 & 45 & 45 & 9/9 \\
Local bound & 27 & 0 & 0 & 9/9 \\
Native AGT & 27 & 0 & 0 & 9/9 \\
\bottomrule\end{tabular}\end{table}
\begin{table}[t]
\centering\small
\caption{Finite model episodes, separated by serving condition. Obs./plan counts
episodes without transport failure or an endpoint stop. Clean gives goal completion;
failure ctx. gives a safe granted effect in that context; reactive requires a recorded
primary-operation error followed by a granted fallback effect. Denominators are observed context episodes; each context originally
has nine planned goals per guard. U is an observed effect count. Missing free-service
episodes are not successful restraint and prevent a complete endpoint comparison.}
\label{tab:model-episodes}
\begin{tabular}{llrrrrr}\toprule
Condition & Guard & Obs./plan & Clean & Failure ctx. & Reactive & U\\\midrule
Local Llama & Allow all & 27/27 & 9/9 & 3/9 & 0/9 & 15 \\
Local Llama & Deny all & 27/27 & 0/9 & 0/9 & 0/9 & 0 \\
Local Llama & Name only & 27/27 & 9/9 & 3/9 & 0/9 & 0 \\
Local Llama & Native AGT & 27/27 & 9/9 & 3/9 & 0/9 & 0 \\
\midrule
Free service & Allow all & 21/27 & 7/7 & 7/7 & 7/7 & 0 \\
Free service & Deny all & 16/27 & 0/6 & 0/2 & 0/2 & 0 \\
Free service & Name only & 15/27 & 6/6 & 5/5 & 5/5 & 0 \\
Free service & Native AGT & 17/27 & 5/5 & 8/8 & 8/8 & 0 \\
\bottomrule\end{tabular}\end{table}


Native AGT and the local bound control have zero unauthorized effects on this finite
matrix and complete all nine clean granted goals. Deny all also has zero such
effects, with zero clean utility. Name-only checking admits 63 unauthorized effects
and the exact-precheck arm admits 45; each occurs with positive name/log consistency.
These examples establish a blind spot in that proxy under the authored faults,
not superiority over the third-party bound baseline: the bound controls tie here.

\subsection{Model-generated actions and operation-failure feedback}
We retain an earlier Qwen3-4B development result of 12 abstentions, including six
explicitly granted tasks. It does not supply efficacy evidence. The new unconstrained
JSON interface also tests development feasibility before evaluation. The Liquid
endpoint emits incompatible native tool-call tokens in all six development episodes.
A subsequently added Cohere candidate has five parse failures, mostly empty content
at a length finish. Neither candidate qualifies for expansion. These additions and
format failures remain exploratory, not an initially fixed model comparison.

Two serving conditions enter separate finite panels: a byte-verified local
Llama3.2-3B GGUF (full hash in the run manifest) and a Nemotron-labelled free endpoint.
Only the local weights and runtime bytes are independently hashable here; the free
endpoint supplies requested/returned model labels, without a verified checkpoint.
These are not two verified model-family replications. Local calls are serialized
on CPU; free calls use four-way scheduling, precluding a latency comparison.

Each panel schedules nine goals, three contexts and four guards, for 108 episodes.
The contexts are a clean granted action, a pending action, and a primary operation
that fails with a separately granted fallback action/ticket. The last is a constructed
environmental fault, while proposals and recovery choices are model generated.
The model sees the complete structured permission, gets actual operation feedback,
and has at most two turns. An exact JSON object, optionally in a JSON fence, is
required; output is not repaired or constrained by a decoding grammar. Requested
serving options are temperature zero, seed \texttt{20261014} and 256 output tokens;
they do not guarantee identical or cap-respecting upstream execution. The declared
CPU timeout is 180 seconds; the free-service timeout is 40 seconds.

The local panel records 180 requests across all 108 episodes. Its statuses are 36 effect-triggered stops (the controller label is completed), 63 model stops, 7 turn limits, 2 parse failures and 0 transport failures.
The free panel attempts 102 requests: 38 episodes complete, 29 stop, and 2 reach the turn limit; 4 have transport errors and 35 are skipped after the endpoint stop.
On the local panel, allow all admits 15 unauthorized effects, all 15 in pending contexts. Name-only and native AGT each admit 0 and 0 such effects, respectively. Native AGT completes 9/9 clean goals. In failure contexts it has 3/9 safe granted effects, all from direct fallback selection; strict failure-then-fallback recovery is 0/9, with 6 observed primary-operation failures.
The original verifier field named safe\_authorized\_completion checks
permission-compatible effects, not compliance with the requested primary-first order.
Direct fallback selection therefore supports a granted effect but not observed
adaptation to failure feedback. A sequence-specific trace check was added on referee
inspection as a descriptive secondary audit, not a prespecified confirmatory endpoint.
Original records and field values remain unchanged. Clean utility has a single exact
requested action; failure-context effects are kept separate from ordered goal completion.
This natural-proposal panel does not distinguish the name-only and
action-bound controls on unauthorized effects. The argument-binding blind spot is
established by the separate constructed matrix, not by claiming a natural substitution
distribution from these episodes.

The gateway wrapped upstream HTTP 429 inside HTTP 502. The client detected the
consecutive-error stopping rule but missed the numeric rate-limit code in the body,
so it did not stop at the first such response. This is a recorded protocol deviation;
no failed request was retried and no model was substituted into the missing cells.
Observed absence of unauthorized effects cannot characterize skipped episodes.
The unequal availability also prevents a complete model-condition comparison.

The tables separate clean utility, granted effects in failure contexts, strict reactive recovery, observed unauthorized effects
and availability. Pending tasks have no granted completion target. Model stops,
parse failures and transport errors are never useful completion. Repeated contexts,
guards and turns of one base goal are not independent tasks. Because these are
small, authored panels with adaptive feasibility selection and partial upstream
availability, we report finite counts without population tests or confidence intervals.
